\documentclass[10pt]{article}

\usepackage[utf8]{inputenc}

\usepackage[T1]{fontenc}

\usepackage{amsmath}

\usepackage{amsfonts}

\usepackage{amssymb}

\usepackage[version=4]{mhchem}

\usepackage{stmaryrd}



\DeclareUnicodeCharacter{0131}{$\imath$}



\begin{document}

Сумма всех малых П. совпадает с пересечением всех максимальных П. Левый идеал $I$ принадлежит радикалу Джекобсона тогда и только тогда, когда IM мал в $M$ для всякого конечно порожденного левого модуля $\boldsymbol{M}$. Элементы малого П. являются н е о б р аз у ю щ и м и, т. е. любая система образующих модуля остается таковой после удаления любого из этих элементов (это, конечно, не означает, что их можно удалить все сразу!). Радикал Джекобсона кольца эндоморфизмов модуля совіадает с множеством эндоморфизмов, имеющих малый образ.



Лит.. [1] К а ш Ф, Модули и кольца, пер. с нем., М., 1981, [2] $\Phi$ е й с К., Алгебра кольца, модули и категории, пер. с англ., т. 1-2, М., 1977-79.



ПоДоБИЕ - преобразование евклидова пространства, при к-ром для любых двух точек $A, B$ и их образов $A^{\prime}, B^{\prime}$ имеет место соотношение $\left|A^{\prime} B^{\prime}\right|=k|A B|$, где $k$ - положительное число, называемое к о эффи ц и е п т о м П.



Каждая гомотетия является подобием. Каждое движение (в том числе и тождественное) также можно рассматривать как преобразование П. с коәффициентом $k$, равным единице. Фигура $F$ наз. пा о д о б н о й фигуре $F^{\prime}$, если существует преобразование П., при к-ром $F \rightarrow F^{\prime}$. П. фигур является отношением эквивалентности, т. е. облладает свойствами рефлексивности, симметричности и транзитивности. П. есть взаимно однозначное отображение евклидова пространства ва себя; П. сохраняет порядок точек на прямой, т. е. если точка $B$ лежит между точками $A, C$ и $B^{\prime}, A^{\prime}$, $C^{\prime}$ - соответствующие их образы при нек-ром П., то $B^{\prime}$ также лежит между точками $A^{\prime}$ и $C^{\prime} ;$ точки, не лежащие на прямой, при любом П. переходят в точки, не лежашие на одной прямой. П. преобразует прямую в прямую, отрезок в отрезок, луч в луч, угол в угол, окружность в окружность. При П. угол сохраняет величину.



II. с коэффициентом $k \neq 1$, преобразующее каждую прямую в параллельную ей прямую, является гомотетией с коэффициентом $k$ или $-k$. Каждое II. можно рассматривать как комІозицию движения $D$ и нек-рой гомотетии $\Gamma$ с положительным коэффициентом.



П. наз. с об с т в е н ны м (не с об с т в е нн ы м), если движение $D$ является собственным (несобственным). Собственное П. сохраняет ориентацию фигур, а несобственное - изменяет ориентацию на противоположную.



Аналогично определяется П. (с сохранением указанных выше свойств) в 3-мерном евклидовом пространстве, а также в $n$-мерном евклидовом и псевдоевклидовом пространствах.



В $n$-мерных римановых, псевдоримановых и финслеровых пространствах П. определяется как преобразование, переводящее метрику пространства в себя с точностью до постоянного множителя.



Совокупность всех П. $n$-мерного евклидова, псевдоевклидова, риманова, псевдориманова или финслерова пространства составляет $r$-членную группу преобразований Ли, наз. г р у п по й п од о б н ы х (г о м о т е т и ч е с к и х) преобразований соответствующего пространства. В каждом из пространств указанных типов $r$-членная группа подобных преобразований Ли содержит (r-1)-членную нормальную подгруппу движений.



Представляют интерес метрич. пространства векторных плотностей с группами П. и изометрий, содержащими бесконечномерные подгруппы изометрий с общими траекториями.



ПОДОБИЯ ТЕОРИЯ - учение об исследовании физич. явлений, основанное на понятии о физич. подобии.



Два физич. явления пा о д о б н ы, если по численным значениям характеристик одного явления можно получить численные значения характеристик другого явления простым пересчетом, к-рый аналогичен переходу от одной системы единиц измерения к другой. Для всякой совокупности подобных явлений все соответствующие безразмерные характеристики (безразмерные комбинации из размерных величин) имеют одинаковое численное значение (см. Размерностей ана$л и з)$. Обратное заключение тоже верно, т. е. если все соответствующие безразмерные характеристики для двух явлений одинаковы, то эти явления физічески подобны.



Анализ размерностей и П. т. тесно связаны между собой и положены в основу экспериментов с моделями. В таких экспериментах осуществляются замены изучения нек-рого явления в натуре изучением аналогичного явления на модели меньшего или большего масштаба (обычно в специальных лабораторных условиях).



После установления системы параметров, определяющих выделенный класс явлений, устанавливаются условия подобия двух явлений. Именно, пусть явление определяется $n$ независимыми параметрами, нек-рые из к-рых могут б́ыть безразмерными. Пусть, далее, размерности определяющих переменных и физич. постоянных выражены через размерности $k$ из этих параметров с независимыми размерностями $(k \leqslant n)$. Тогда из $n$ величин можно составить только $n-k$ независимых безразмерных комбинаций. Все искомые безразмерные характеристики явления можно рассматривать как функции от этих $n-k$ независимых безразмерных комбинаций, составленных из определяющих параметров. Среди всех безразмерных величин, составленных из определяющих характеристик явления, всегда можно указать нек-рую базу, т. е. систему безразмерных величин, к-рые определяют собой все остальные.



Определенный соответствующей постановкой задачи класс явлений содержит явления, вообще неподобные между собой. Выделение из него подкласса подобных явлений осуществляется с помощью следующего условия.



Для подобия двух явлений необходимо и достаточно, чтобы численные значения безразмерных комбинаций, составленных из полного перечня определяющих параметров, образующих базу, в этих двух явлениях были одинаковы. Условия о постоянстве базы́ отвлеченных параметров, составленных из заданных величин, определяющих явление, наз. к р и т е р и я м и по д об и я.



В гидродинамике важнейшими критериями подобия являются Рейнольдса число, характеризующее соотношение между инерционными силами и силами вязкости, Маха число, учитывающее сжимаемость газа, и Фруда число, характеризующее соотношение между инерционными силами и силами тяжести. Основными критериями подобия процессов теплопередачи между жидкостью (газом) и обтекаемым телом являются: Прандтля число, характеризующее термодинамич. состояние среды; Нуссельта число, характеризующее интенсивность конвективного теплообмена между поверхностыю тела и потоком жидкости (газа); Пекле число, характеризующее соотнопение между конвективным и молекулярным процессами переноса тепла в жидкости; Стәнтона число, характеризующее интенсивность диссипации энергии в потоке жидкости или газа. Для распределения тепла в твердом теле критериями подобия являются Фурье число, характөризующее скорость изменения тепловых условий в окружающей среде и скорость перестройки поля темп-ры внутри тела, и число Био, определяющее характер соответствия между температурными условиями среды и расıределением температуры внутри тела. В про-





\end{document}